\documentclass[11pt,a4paper]{article}

\usepackage[T1]{fontenc}
\usepackage[utf8]{inputenc}
\usepackage[french]{babel}
\usepackage{amsmath,amssymb,amsthm}
\usepackage{mathtools}
\usepackage[margin=2.6cm]{geometry}
\usepackage{enumitem}
\usepackage{xcolor}
\usepackage[colorlinks=true, linkcolor=blue!60!black, citecolor=blue!60!black]{hyperref}

% ----- Environnements de théorèmes -----
\theoremstyle{plain}
\newtheorem{theorem}{Théorème}[section]
\newtheorem{proposition}[theorem]{Proposition}
\newtheorem{lemma}[theorem]{Lemme}
\newtheorem{corollary}[theorem]{Corollaire}
\theoremstyle{definition}
\newtheorem{example}[theorem]{Exemple}
\theoremstyle{remark}
\newtheorem{remark}[theorem]{Remarque}

% ----- Macros -----
\newcommand{\E}{\mathbb{E}}
\newcommand{\Prob}{\mathbb{P}}
\newcommand{\R}{\mathbb{R}}
\newcommand{\Q}{\mathbb{Q}}
\newcommand{\N}{\mathbb{N}}
\newcommand{\F}{\mathcal{F}}
\newcommand{\ps}{\text{p.s.}}
\DeclarePairedDelimiter{\abs}{\lvert}{\rvert}
\DeclarePairedDelimiter{\norm}{\lVert}{\rVert}

\title{Convergence des martingales de Doob\\et dichotomie de Kakutani}
\author{}
\date{Juillet 2026}

\begin{document}
\maketitle

\begin{abstract}
Ce document présente le théorème de convergence presque sûre des martingales de Doob,
l'inégalité des montées qui en constitue le c\oe ur, puis la dichotomie de Kakutani
(1948) pour les martingales produit et les mesures produit infinies, avec sa preuve
complète. On détaille le cas gaussien, sa reformulation comme théorème de
Cameron--Martin, et les liens avec la statistique asymptotique et la finance
mathématique.
\end{abstract}

\tableofcontents

%=======================================================================
\section{Le théorème de convergence de Doob}
%=======================================================================

Dans toute la suite, $(\Omega, \F, (\F_n)_{n \ge 0}, \Prob)$ désigne un espace de
probabilité filtré.

\begin{theorem}[Doob]\label{thm:doob}
Soit $(X_n)_{n\ge 0}$ une sous-martingale telle que
\begin{equation}\label{eq:hyp}
\sup_{n\ge 0} \E\bigl[X_n^+\bigr] < \infty.
\end{equation}
Alors il existe une variable aléatoire $X_\infty$ telle que
\[
X_n \xrightarrow[n\to\infty]{\ps} X_\infty,
\qquad \E\bigl[\abs{X_\infty}\bigr] < \infty.
\]
\end{theorem}

\begin{remark}
Pour une sous-martingale, la condition \eqref{eq:hyp} équivaut à la bornitude dans
$L^1$ : en effet $\E[\abs{X_n}] = 2\,\E[X_n^+] - \E[X_n]$ et la suite
$(\E[X_n])_n$ est croissante, donc minorée par $\E[X_0]$.
\end{remark}

\begin{corollary}\label{cor:cas-part}
Les processus suivants convergent presque sûrement :
\begin{enumerate}[label=(\roman*)]
\item toute martingale bornée dans $L^1$ ;
\item toute sous-martingale majorée par une constante ;
\item toute surmartingale positive --- sans aucune hypothèse d'intégrabilité
      supplémentaire, puisque $\E[X_n^+] = \E[X_n] \le \E[X_0]$.
\end{enumerate}
\end{corollary}

\subsection{L'inégalité des montées}

Pour $a < b$ réels, on note $U_n[a,b]$ le nombre de \emph{traversées montantes}
complètes de l'intervalle $[a,b]$ par la trajectoire $(X_0, \dots, X_n)$ :
le nombre de fois où le processus, parti sous le niveau $a$, atteint le niveau $b$.

\begin{lemma}[Inégalité des montées de Doob]\label{lem:upcrossing}
Si $(X_n)$ est une sous-martingale, alors pour tous $a<b$ et tout $n \ge 0$,
\[
(b-a)\, \E\bigl[U_n[a,b]\bigr]
\;\le\;
\E\bigl[(X_n - a)^+\bigr] - \E\bigl[(X_0 - a)^+\bigr].
\]
\end{lemma}

\begin{proof}[Preuve du théorème~\ref{thm:doob} à partir du lemme]
Sous l'hypothèse \eqref{eq:hyp}, le membre de droite est majoré par
$\sup_n \E[X_n^+] + \abs{a} < \infty$ uniformément en $n$. Par convergence monotone,
$\E[U_\infty[a,b]] < \infty$, donc $U_\infty[a,b] < \infty$ presque sûrement.
Or
\[
\Bigl\{\liminf_n X_n < \limsup_n X_n\Bigr\}
= \bigcup_{\substack{a<b \\ a,b \in \Q}} \bigl\{U_\infty[a,b] = \infty\bigr\},
\]
union dénombrable d'événements négligeables. Donc $X_n$ converge p.s.\ dans
$\overline{\R}$ vers une limite $X_\infty$. Le lemme de Fatou donne enfin
$\E[\abs{X_\infty}] \le \liminf_n \E[\abs{X_n}] < \infty$, ce qui assure en
particulier $\abs{X_\infty} < \infty$ p.s.
\end{proof}

\subsection{Convergence $L^1$ et martingales fermées}

La convergence du théorème~\ref{thm:doob} a lieu presque sûrement, mais pas
nécessairement dans $L^1$.

\begin{proposition}\label{prop:UI}
Soit $(X_n)$ une martingale. Les assertions suivantes sont équivalentes :
\begin{enumerate}[label=(\roman*)]
\item la famille $(X_n)_{n \ge 0}$ est uniformément intégrable ;
\item $X_n$ converge dans $L^1$ ;
\item la martingale est \emph{fermée} : il existe $Z \in L^1$ telle que
      $X_n = \E[Z \mid \F_n]$ pour tout $n$.
\end{enumerate}
Dans ce cas $X_n \to X_\infty$ p.s.\ et dans $L^1$, et
$X_n = \E[X_\infty \mid \F_n]$.
\end{proposition}

La dichotomie de Kakutani, objet de la section suivante, fournit un critère
\emph{exact} d'intégrabilité uniforme pour une classe importante de martingales :
les martingales produit.

%=======================================================================
\section{La dichotomie de Kakutani}
%=======================================================================

\subsection{Formulation martingale}

\begin{theorem}[Kakutani, 1948 --- version martingale]\label{thm:kakutani-mart}
Soit $(\xi_k)_{k \ge 1}$ une suite de variables aléatoires indépendantes, positives,
d'espérance $\E[\xi_k] = 1$, et soit
\[
X_n = \prod_{k=1}^{n} \xi_k, \qquad X_0 = 1,
\]
la martingale produit associée. On pose $a_k = \E\bigl[\sqrt{\xi_k}\bigr] \in (0,1]$.
Alors exactement l'une des deux situations suivantes se produit :
\begin{enumerate}[label=(\Alph*)]
\item $\displaystyle\prod_{k\ge 1} a_k > 0$
      \emph{(de manière équivalente $\sum_k (1-a_k) < \infty$)} :
      la martingale $(X_n)$ est uniformément intégrable, $X_n \to X_\infty$
      p.s.\ et dans $L^1$, et $\E[X_\infty] = 1$ ;
\item $\displaystyle\prod_{k\ge 1} a_k = 0$ :
      \quad $X_\infty = 0$ presque sûrement.
\end{enumerate}
\end{theorem}

\begin{remark}
Noter que $a_k \le 1$ résulte de l'inégalité de Cauchy--Schwarz
($\E[\sqrt{\xi_k}] \le \sqrt{\E[\xi_k]} = 1$), avec égalité si et seulement si
$\xi_k = 1$ p.s. Le produit $\prod_k a_k$ mesure donc l'accumulation des
perturbations : chaque facteur non trivial coûte strictement.
\end{remark}

\begin{proof}
Dans les deux cas, $(X_n)$ est une martingale positive, donc converge p.s.\ vers
une limite $X_\infty \ge 0$ intégrable (corollaire~\ref{cor:cas-part}).

\medskip
\noindent\emph{Cas (A).} Posons
\[
Y_n = \prod_{k=1}^{n} \frac{\sqrt{\xi_k}}{a_k}.
\]
Par indépendance, $(Y_n)$ est une martingale positive de moyenne $1$, et
\[
\E\bigl[Y_n^2\bigr]
= \prod_{k=1}^{n} \frac{\E[\xi_k]}{a_k^2}
= \prod_{k=1}^{n} a_k^{-2}
\;\le\; \Bigl(\prod_{k\ge 1} a_k\Bigr)^{-2} < \infty.
\]
La martingale $(Y_n)$ est donc bornée dans $L^2$ : elle converge dans $L^2$
(et p.s.) vers une limite $Y_\infty$. Comme
\[
X_n = Y_n^2 \prod_{k=1}^{n} a_k^2,
\]
et que $Y_n^2 \to Y_\infty^2$ dans $L^1$ (convergence $L^2$ de $Y_n$ plus
bornitude $L^2$), la suite $(X_n)$ converge dans $L^1$. La
proposition~\ref{prop:UI} donne l'intégrabilité uniforme et
$\E[X_\infty] = \lim_n \E[X_n] = 1$ ; en particulier
$\Prob(X_\infty > 0) > 0$.

\medskip
\noindent\emph{Cas (B).} Par indépendance,
\[
\E\bigl[\sqrt{X_n}\bigr] = \prod_{k=1}^{n} a_k \xrightarrow[n\to\infty]{} 0.
\]
Comme $\sqrt{X_n} \to \sqrt{X_\infty}$ p.s., le lemme de Fatou donne
\[
\E\bigl[\sqrt{X_\infty}\bigr] \le \liminf_n \E\bigl[\sqrt{X_n}\bigr] = 0,
\]
donc $X_\infty = 0$ p.s.
\end{proof}

\begin{remark}[Tout ou rien]
Il n'existe aucun régime intermédiaire : $\Prob(X_\infty > 0)$ vaut soit une
quantité rendant la martingale fermée avec $\E[X_\infty]=1$, soit exactement $0$.
On peut le relier à une loi du zéro--un : l'événement $\{X_\infty > 0\}$ est,
modulo un ensemble négligeable, mesurable par rapport à la tribu asymptotique des
$(\xi_k)$, donc de probabilité $0$ ou $1$ par la loi de Kolmogorov.
\end{remark}

\subsection{Formulation en théorie de la mesure}

\begin{theorem}[Kakutani --- mesures produit]\label{thm:kakutani-mes}
Soient $(E_k, \mathcal{E}_k)_{k\ge1}$ des espaces mesurables et, sur chacun, deux
probabilités équivalentes $P_k \sim Q_k$. On définit sur l'espace produit
$E = \prod_k E_k$ les mesures
\[
P = \bigotimes_{k\ge1} P_k, \qquad Q = \bigotimes_{k\ge1} Q_k,
\]
et l'\emph{affinité de Hellinger} de chaque facteur
\[
\rho_k = \int_{E_k} \sqrt{\frac{dQ_k}{dP_k}}\; dP_k \;\in\; (0,1].
\]
Alors :
\begin{enumerate}[label=(\Alph*)]
\item si $\prod_k \rho_k > 0$, alors $Q \sim P$ (mesures équivalentes) ;
\item si $\prod_k \rho_k = 0$, alors $Q \perp P$ (mesures étrangères).
\end{enumerate}
\end{theorem}

\begin{proof}[Lien avec la version martingale]
Travaillons sous $P$ et posons $\xi_k = \frac{dQ_k}{dP_k}(\omega_k)$, variables
indépendantes, positives, d'espérance $1$ sous $P$. La martingale produit
$X_n = \prod_{k \le n} \xi_k$ n'est autre que le processus de densité
\[
X_n = \frac{dQ}{dP}\bigg|_{\F_n},
\qquad \F_n = \sigma(\omega_1, \dots, \omega_n),
\]
et $a_k = \E_P[\sqrt{\xi_k}] = \rho_k$. Dans le cas (A), la martingale est fermée
par $X_\infty$ avec $\E_P[X_\infty]=1$ ; on vérifie alors que
$dQ = X_\infty\, dP$ avec de plus $X_\infty > 0$ $P$-p.s.\ (appliquer le même
argument en échangeant les rôles de $P$ et $Q$, licite car
$\int \sqrt{dP_k/dQ_k}\,dQ_k = \rho_k$ également), d'où $Q \sim P$. Dans le cas
(B), $X_\infty = 0$ $P$-p.s.\ tandis que, par symétrie, le processus de densité
inverse tend vers $0$ $Q$-p.s. : les deux mesures se concentrent sur des
événements disjoints, $Q \perp P$.
\end{proof}

\begin{remark}
Le contenu frappant du théorème est l'absence de décomposition de Lebesgue non
triviale : bien que chaque facteur soit équivalent, une infinité de perturbations
individuellement bénignes produit soit une équivalence globale, soit une
singularité totale --- jamais un mélange.
\end{remark}

%=======================================================================
\section{Le cas gaussien et le théorème de Cameron--Martin}
%=======================================================================

\begin{example}[Martingale exponentielle gaussienne]\label{ex:gauss}
Soient $(Z_k)$ i.i.d.\ $\mathcal{N}(0,1)$, $(\sigma_k)$ une suite réelle, et
\[
\xi_k = \exp\Bigl(\sigma_k Z_k - \tfrac{\sigma_k^2}{2}\Bigr),
\qquad
X_n = \prod_{k=1}^n \xi_k
= \exp\Bigl(\sum_{k=1}^n \sigma_k Z_k - \tfrac12 \sum_{k=1}^n \sigma_k^2\Bigr).
\]
La transformée de Laplace gaussienne $\E[e^{\lambda Z}] = e^{\lambda^2/2}$ donne
$\E[\xi_k] = 1$ et
\[
a_k = \E\bigl[\sqrt{\xi_k}\bigr]
= \E\Bigl[e^{\frac{\sigma_k}{2} Z_k - \frac{\sigma_k^2}{4}}\Bigr]
= e^{\frac{\sigma_k^2}{8} - \frac{\sigma_k^2}{4}}
= e^{-\sigma_k^2/8}.
\]
La dichotomie se lit alors sur une condition de sommabilité :
\[
\prod_k a_k > 0
\iff
\boxed{\;\sum_{k\ge1} \sigma_k^2 < \infty.\;}
\]
\begin{itemize}
\item Si $\sum_k \sigma_k^2 < \infty$ : $X_n \to X_\infty > 0$ p.s.,
      $\E[X_\infty]=1$, la martingale est fermée.
\item Si $\sum_k \sigma_k^2 = \infty$ (en particulier à $\sigma_k = \sigma > 0$
      constant) : $X_\infty = 0$ p.s., alors que $\E[X_n] = 1$ pour tout $n$.
      À $\sigma$ constant, la loi des grands nombres précise la vitesse :
      $\frac1n \log X_n \to -\sigma^2/2$ p.s.
\end{itemize}
\end{example}

\begin{corollary}[Cameron--Martin, cas discret]
Sur $\R^{\N}$ muni des tribus produit,
\[
\bigotimes_{k\ge1} \mathcal{N}(\theta_k, 1) \;\sim\; \bigotimes_{k\ge1} \mathcal{N}(0,1)
\iff
(\theta_k)_k \in \ell^2,
\]
et il y a singularité totale sinon.
\end{corollary}

\begin{proof}
La densité du facteur $k$ est
$\frac{d\mathcal{N}(\theta_k,1)}{d\mathcal{N}(0,1)}(x) = e^{\theta_k x - \theta_k^2/2}$,
soit exactement $\xi_k$ de l'exemple~\ref{ex:gauss} avec $\sigma_k = \theta_k$.
La condition $\sum \theta_k^2 < \infty$ conclut par le
théorème~\ref{thm:kakutani-mes}.
\end{proof}

\begin{remark}
C'est la version discrète du théorème de Cameron--Martin pour la mesure de
Wiener : une translation du brownien par une fonction $h$ produit une mesure
équivalente si et seulement si $h$ appartient à l'espace de Cameron--Martin
(dérivée dans $L^2$), et une mesure étrangère sinon --- même mécanisme, même
dichotomie.
\end{remark}

%=======================================================================
\section{Portée du résultat}
%=======================================================================

\paragraph{Statistique asymptotique.}
Soit à tester $H_0 : \text{les observations suivent } P$ contre
$H_1 : \text{elles suivent } Q$, à partir de la suite infinie d'observations
indépendantes. La dichotomie affirme que la discrimination parfaite est soit
possible (cas singulier : le test $\{X_\infty = 0\}$ ne commet asymptotiquement
aucune erreur), soit impossible même asymptotiquement (cas équivalent : aucun test
consistant n'existe, les deux hypothèses restent \emph{contiguës}). L'affinité de
Hellinger $\rho_k$, et la distance de Hellinger associée
$H^2(P_k, Q_k) = 2(1 - \rho_k)$, sont ainsi les quantités naturelles de la théorie
de la contiguïté de Le~Cam : la frontière $\sum_k H^2(P_k, Q_k) < \infty$ sépare
exactement le détectable de l'indétectable.

\paragraph{Finance mathématique.}
Le processus de densité d'un changement de mesure de Girsanov,
$X_t = \exp\bigl(\int_0^t \theta_s\, dW_s - \frac12 \int_0^t \theta_s^2\, ds\bigr)$,
est une martingale exponentielle du type de l'exemple~\ref{ex:gauss}. Sur horizon
fini, la condition de Novikov
$\E\bigl[e^{\frac12\int_0^T \theta_s^2 ds}\bigr] < \infty$ garantit qu'elle est une
vraie martingale fermée et que la mesure risque-neutre existe. Sur horizon infini,
c'est le mécanisme de Kakutani qui gouverne : si la \og prime de risque \fg{}
cumulée $\int_0^\infty \theta_s^2\, ds$ diverge, le changement de mesure dégénère
($X_\infty = 0$), et les mesures historique et risque-neutre deviennent étrangères
--- même mathématique que le cas discret, avec $\sum \sigma_k^2$ remplacée par
l'intégrale.

\paragraph{En résumé.}
Le théorème de Doob garantit l'existence de la limite ; la dichotomie de Kakutani
répond à la question que Doob laisse ouverte --- \emph{la limite conserve-t-elle la
masse ?} --- et y répond de la manière la plus tranchée possible : tout ou rien,
avec un critère exact, calculable facteur par facteur.

\begin{thebibliography}{9}
\bibitem{kakutani} S.~Kakutani, \emph{On equivalence of infinite product measures},
Annals of Mathematics \textbf{49} (1948), 214--224.
\bibitem{williams} D.~Williams, \emph{Probability with Martingales},
Cambridge University Press, 1991.
\bibitem{durrett} R.~Durrett, \emph{Probability: Theory and Examples},
5\textsuperscript{e} éd., Cambridge University Press, 2019.
\bibitem{lecam} L.~Le~Cam, \emph{Asymptotic Methods in Statistical Decision Theory},
Springer, 1986.
\end{thebibliography}

\end{document}
